\documentclass[11pt]{amsart}
\usepackage[T1]{fontenc}
\usepackage{lmodern,microtype,mathtools,amssymb,booktabs,array}
\usepackage[margin=30mm]{geometry}
\usepackage[colorlinks=true,linkcolor=blue,citecolor=blue,urlcolor=blue]{hyperref}
\hypersetup{pdftitle={Betti numbers of numerical semigroup tangent cones and interval completion},
pdfauthor={Keisuke Sato},pdfkeywords={numerical semigroup, tangent cone, Betti numbers, interval completion}}
\newtheorem{theorem}{Theorem}[section]
\newtheorem{proposition}[theorem]{Proposition}
\newtheorem{lemma}[theorem]{Lemma}
\newtheorem{corollary}[theorem]{Corollary}
\theoremstyle{definition}
\newtheorem{definition}[theorem]{Definition}
\theoremstyle{remark}
\newtheorem{remark}[theorem]{Remark}
\newcommand{\N}{\mathbb N}
\newcommand{\gr}{\operatorname{gr}}
\newcommand{\Tor}{\operatorname{Tor}}
\newcommand{\ord}{\operatorname{ord}}
\newcommand{\Soc}{\operatorname{Soc}}
\newcommand{\Ap}{\operatorname{Ap}}
\newcommand{\reg}{\operatorname{reg}}
\DeclareMathOperator{\im}{im}
\title[Betti numbers and interval completion]{Betti numbers of numerical semigroup tangent cones\texorpdfstring{\\}{ }and interval completion}
\author[Keisuke Sato]{Keisuke Sato}
\address{Independent researcher, Japan}
\email{stksk310@gmail.com}
\urladdr{https://orcid.org/0009-0005-4230-0268}
\subjclass[2020]{Primary 13D02; Secondary 13F65, 20M14}
\keywords{Numerical semigroup, tangent cone, Betti numbers, interval completion,
formal verification}
\date{}
\begin{document}
\raggedbottom
\begin{abstract}
Let $S$ be a numerical semigroup with smallest and largest minimal
generators $m$ and $M$, and let $\widetilde S$ be the semigroup generated
by all integers in $[m,M]$. We prove that, over every field, each total
Betti number of the tangent cone of $k[[t^S]]$ is at most the corresponding
Betti number for $k[[t^{\widetilde S}]]$. Each tangent cone is presented
over its own minimal polynomial ring. The proof combines an Ap\'ery
filtration for ordinary interval completions with an endpoint reduction
for narrow intervals. The latter retains the torsion contribution and
therefore applies without a Cohen--Macaulay hypothesis. A Lean
formalization establishes the comparison for every numerical semigroup,
every field, and every homological degree. As a consequence, we obtain
the sharp width bound $\beta_i(k[[t^S]];k)\le i\binom{w+1}{i+1}$
for $i\ge1$, where $w=M-m$.
\end{abstract}
\maketitle
\section{Introduction}
\label{sec:introduction}
Let $S=\langle m=g_1<g_2<\cdots<g_e=M\rangle$ be a numerical
semigroup, written with its minimal generators. Over a field $k$, its
completed semigroup ring has tangent cone
\[
 G_S=\gr_{\mathfrak m} k[[t^S]],\qquad
 \mathfrak m=(t^{g_1},\ldots,t^{g_e}).
\]
The initial forms of the minimal generators give a surjection
$P_S=k[X_1,\ldots,X_e]\twoheadrightarrow G_S$.
We write
\[
 \beta_i(G_S;k)=\dim_k\Tor_i^{P_S}(G_S,k)
\]
for its total Betti numbers. Thus $\beta_0=1$, and the positive rows
measure the numbers of generators and higher syzygies of the defining
ideal, with a shift of one in the homological index.

The interval completion of $S$ is
\[
 \widetilde S=\langle m,m+1,\ldots,M\rangle.
\]
Its listed generators need not be minimal. Throughout, Betti numbers
on the right are computed over the polynomial ring on the minimal
generators of $\widetilde S$, independently of the presentation of $S$.

\begin{theorem}\label{thm:main}
For every numerical semigroup $S$, every field $k$, and every $i\ge0$,
\[
 \beta_i(G_S;k)\le\beta_i(G_{\widetilde S};k).
\]
\end{theorem}

No restriction is imposed on the characteristic of $k$ or on the
Cohen--Macaulayness of $G_S$. The use of different minimal polynomial
rings is essential: the interval completion may have many more minimal
generators than the original semigroup. The theorem concerns the total
Betti numbers; it does not assert a comparison in each internal degree.

Theorem~\ref{thm:main} resolves Problem~14 in the problem list of
Moscariello--Sammartano \cite{ProblemList}. Herzog--Stamate had
formulated the first-row comparison and discussed its higher-syzygy
extension \cite{APCone}; the problem list states the all-row form used
here.

The target has an explicit Betti table. Put $w=M-m$. If $m\ge2$, its
minimal generators are $m,m+1,\ldots,m+v$, where
$v=\min\{w,m-1\}$. The arithmetic-sequence formula of
Gimenez--Sengupta--Srinivasan, together with the tangent-cone comparison
of Herzog--Stamate, determines its Betti numbers
\cite{APFormula,APCone}. We use that formula as a numerical capacity
against which to compare the source.

The comparison also yields a sharp bound in terms of width alone,
settling Conjecture~1.3 of Caviglia--Moscariello--Sammartano
\cite{CMSBounds}. For the complete local ring
$R_S=k[[t^S]]$, let $\beta_i(R_S;k)$ denote its Betti numbers over
the minimal regular presentation $k[[X_1,\ldots,X_e]]$.
\begin{corollary}\label{cor:width-bound}
For every numerical semigroup $S$ of width $w$, every field $k$, and
every $i\ge1$,
\[
 \beta_i(R_S;k)\le\beta_i(G_S;k)
                 \le i\binom{w+1}{i+1}.
\]
For each $w\ge1$ these bounds are attained simultaneously by
$S=\langle w+1,w+2,\ldots,2w+1\rangle$.
\end{corollary}
The proof, including the local-ring comparison, is given at the end of
Section~\ref{sec:preliminaries}.

The proof separates at $w=m-1$. For $w\ge m-1$, the completion is
the ordinary semigroup $\{0\}\cup[m,\infty)$, and its $i$th positive
Betti number is $i\binom{m}{i+1}$. An Ap\'ery quotient of length $m$
organizes both the regular part and the torsion of $G_S$. One sequence
of linear cuts yields a single resource profile controlling all rows.
The estimates are numerical after that construction; low cut ranks and
large cut ranks require different inequalities.

For $w\le m-2$, divide the offsets $g_j-m$ by their common divisor
$d$. This gives a normalized width $u=w/d$, but leaves the original
semigroup and its physical width $du$ unchanged. Keeping the two endpoint
generators separates the proof into a finite internal core and an
endpoint kernel. The kernel is determined by all competing
representations, and its socle supplies the correction to the core Betti
bound. Estimates for finite cyclic groups and lex quotients then close
the wide normalized ranges. The remaining cases involve one or two
internal variables and are handled directly.

Lan--Tu--Vu proved the interval-completion comparison for three
generators \cite{LanTuVu}. The proof below includes a compatible
one-internal-variable argument, so that the small-generator boundary
uses the same mechanism as the general narrow case. Our use of Green's
restriction theorem, lex Betti bounds, Kneser's theorem, and the
Savchev--Chen theorem is stated where each enters the argument.

The accompanying Lean development verifies the theorem for the actual
completed tangent cones and their minimal presentations. The mathematical
proof is given in Sections~\ref{sec:preliminaries}--\ref{sec:synthesis};
Appendix~\ref{sec:formalization} identifies the formal statement and the
source version used here.

\subsection*{AI Declaration}
Generative AI was used substantially in the mathematical research and
preparation of this manuscript. OpenAI ChatGPT 6 Astra Pro was used to
explore and construct candidate proofs. OpenAI GPT-6, with the
\emph{extra high} reasoning setting, was used to audit the proposed
arguments, look for errors, and check exceptional cases. OpenAI ChatGPT
Work, using GPT-6.1 Sol with the \emph{Ultra} reasoning setting, was
used for later-stage mathematical proof construction, refinement,
integration, and audit. OpenAI Codex
with GPT-6.1 Sol, using the \emph{high} and \emph{Ultra} reasoning
settings, was used to refine and simplify the proofs, develop the Lean
formalization, draft the English manuscript, and revise its exposition.
The model and setting names identify the interfaces used during the work.

The author's contribution was supervisory: he selected the problem,
designed and supervised an iterative workflow of AI proof construction
and separate AI audit, and directed the organization, simplification,
and presentation of the resulting arguments. The AI contributions
included mathematical reasoning and proof development as well as
writing. Audit rounds within this workflow are AI-assisted checks,
rather than independent human referee reports. The author is responsible
for the content and accuracy of the manuscript. The role of Lean and
the verification code is described separately in
Appendix~\ref{sec:formalization}.

\section{Common reductions and the interval target}
\label{sec:preliminaries}
We use $\binom ab=0$ when $b<0$ or $b>a$, for $a\ge0$.
Unless expressly indicated otherwise, differences in numerical
inequalities are ordinary integer or rational differences. Rational
weights in the proof are counting coefficients, and impose no condition
on the characteristic of $k$.

\subsection{Orders, completion, and scalar extension}
For $s\in S$, let $\ord(s)$ be the largest length of an expression of
$s$ as a sum of minimal generators. If $m$ is the multiplicity, the
Ap\'ery set $\Ap(S,m)=\{s\in S:s-m\notin S\}$ contains one element
in each residue class modulo $m$. Every $s\in S$ has a unique expression
$s=b+\ell m$ with $b\in\Ap(S,m)$ and $\ell\ge0$.

\begin{proposition}\label{prop:bridge}
The following operations preserve the objects needed in
Theorem~\ref{thm:main}.
\begin{enumerate}
\item The associated graded rings of $k[t^S]$ and $k[[t^S]]$ at their
positive monomial ideals are naturally isomorphic, preserving the
initial form of every minimal generator.
\item For every field extension $K/k$,
$\beta_i(G_S;k)=\beta_i(G_S\otimes_k K;K)$ for all $i$.
\item If $w\le m-2$ and $e\ge3$, there are integers $d\ge1$,
$u\ge e-1$, and $0<a_1<\cdots<a_q<u$, with $q=e-2$, such that
\begin{gather}\label{eq:normalization}
 S=\langle m,m+da_1,\ldots,m+da_q,m+du\rangle,\\
 \gcd(a_1,\ldots,a_q,u)=\gcd(m,d)=1,\qquad m\ge du+2.\nonumber
\end{gather}
Here $du=w$ is the physical width.
\end{enumerate}
\end{proposition}
\begin{proof}
The $n$th power of the positive monomial ideal consists exactly of the
elements supported on exponents $s$ with $\ord(s)\ge n$. This holds
also for formal series: assign each such exponent to one of the finitely
many length-$n$ products that divides it, and group its coefficient into
the corresponding formal series. The degree-$n$ quotient consequently
has basis the exponents of order exactly $n$. This set is finite, since
$s=b+\ell m$ of order $n$ has $\ell\le n$. Both associated graded
rings have these bases, and multiplication agrees: a product survives
exactly when orders add. This proves the first assertion, including
compatibility with the polynomial-ring action.

Tensor a minimal graded free resolution with $K$. Flatness preserves
exactness, and all differential entries remain in the homogeneous
maximal ideal, so minimality and ranks are preserved. The preceding
monomial bases also identify the scalar extension of the tangent cone.

For the last assertion, take
$d=\gcd(g_2-m,\ldots,g_e-m)$ and $u=(g_e-m)/d$.
Dividing the internal offsets by $d$ gives the $a_j$. Their common gcd
with $u$ is one, while $\gcd(m,d)=1$ follows from $\gcd(S)=1$.
There are $q$ distinct internal positive integers below $u$, giving
$u\ge q+1=e-1$. The narrow-width hypothesis gives $m\ge du+2$.
\end{proof}

In particular, an infinite extension may be used for a general linear
cut and then removed. We never assume that the whole ring $G_S$ is
finite dimensional over $k$; finiteness is used for its graded pieces,
Artinian quotients, and Koszul homology.

\subsection{Exact target and uniform capacity}
Suppose $m\ge2$, and put
\[
 v=\min\{M-m,m-1\},\qquad
 D=v-((m-1)\bmod v),\quad 1\le D\le v.
\]
The minimal generators of $\widetilde S$ are precisely
$m,m+1,\ldots,m+v$. If $M<2m$, all listed generators are
indecomposable. Otherwise $m,\ldots,2m-1$ generate every integer
at least $m$, by induction after subtraction of $m$; generators above
$2m-1$ are redundant.

\begin{proposition}\label{prop:capacity}
For every $i\ge1$, the target Betti number equals
\begin{equation}\label{eq:capacity}
 F_i(m,v)=i\binom v{i+1}+
 \begin{cases}
 (D+1-i)\binom v{i-1},&i\le D,\\
 (i-D)\binom vi,&i>D.
 \end{cases}
\end{equation}
For $1\le i\le v$, it satisfies
\begin{equation}\label{eq:mincapacity}
 F_i(m,v)\ge F_i^{\min}(v):=
 i\binom v{i+1}+\min\left\{\binom v{i-1},\binom vi\right\}.
\end{equation}
For fixed $i$, $F_i^{\min}(v)$ is nondecreasing on $v\ge i$.
If $v=m-1$, then
\begin{equation}\label{eq:ordinary-capacity}
 F_i(m,m-1)=i\binom m{i+1}.
\end{equation}
\end{proposition}
\begin{proof}
The arithmetic-sequence Betti formula \cite[Theorem~4.1]{APFormula}
and equality of semigroup-ring and tangent-cone Betti numbers
\cite[Proposition~2.5]{APCone} apply to the minimal interval generators.
Writing $m=av+b$ with $1\le b\le v$ gives $D=v-b+1$ and yields
\eqref{eq:capacity}, with the quotient-ring homological indexing.
On the branch $D\ge i$ the correction is at least $\binom v{i-1}$;
on $D<i$ it is at least $\binom vi$. This proves
\eqref{eq:mincapacity}, including $i=1$ and $i=v=1$.
Each binomial coefficient increases with $v$, proving monotonicity.
For $v=m-1$ one has $D=v$; Pascal's identity gives
\eqref{eq:ordinary-capacity}.
\end{proof}

Some useful consequences, with their necessary width ranges, are
\begin{align}\label{eq:low-capacities}
 F_1(m,v)&\ge\binom v2+1 &&(v\ge1),\nonumber\\
 F_2(m,v)&\ge2\binom v3+v &&(v\ge3),\\
 F_3(m,v)&\ge3\binom v4+\binom v2 &&(v\ge5),\nonumber\\
 F_4(m,v)&\ge4\binom v5 &&(v\ge4).\nonumber
\end{align}
We do not extend the third-row estimate to $v=4$.
In the narrow case $v=du\ge u$. Consequently, whenever $u\ge i$,
an upper bound $\beta_i(G_S;k)\le F_i^{\min}(u)$ pays for the
actual target by monotonicity. At the smallest normalized widths, it is
sometimes more efficient to use the physical width directly or to
identify source and target.

\subsection{The width bound for the local ring}
\begin{proof}[Proof of Corollary~\ref{cor:width-bound}]
Put $\widehat P_S=k[[X_1,\ldots,X_e]]$. The standard filtered
Koszul comparison gives
$\beta_i(R_S;k)\le\beta_i(G_S;k)$. Indeed, filter the Koszul
complex on the images of $X_1,\ldots,X_e$ by powers of the maximal
ideal, shifting the filtration by the exterior degree. Its associated
graded complex is the Koszul complex of the initial forms on $G_S$.
The resulting convergent spectral sequence starts with
$\Tor^{P_S}(G_S,k)$ and abuts to the associated graded of
$\Tor^{\widehat P_S}(R_S,k)$. Completeness and the finite-dimensional
first-page homology give the dimension inequality in each homological
degree. This argument works over every field.

If $w=0$, then $S=\N$ and all positive Betti numbers vanish.
Otherwise use Theorem~\ref{thm:main} and put
$v=\min\{w,m-1\}$. In \eqref{eq:capacity}, the correction for
$i\le D$ is at most
\[
 (v+1-i)\binom v{i-1}=i\binom vi;
\]
for $i>D$ it is also at most $i\binom vi$. Hence
\[
 \beta_i(G_{\widetilde S};k)=F_i(m,v)
 \le i\binom{v+1}{i+1}\le i\binom{w+1}{i+1}.
\]
For the stated ordinary semigroup, $m=w+1$ and
\eqref{eq:ordinary-capacity} is equality in the width bound.
The arithmetic-sequence semigroup ring and its tangent cone have the
same total Betti numbers \cite{APFormula,APCone}, giving equality also
for $R_S$.
\end{proof}

\section{The ordinary target}
\label{sec:ordinary}

Throughout this section, let $S=\langle m=g_1<\cdots<g_e\rangle$
be minimally generated, with $e\ge3$, and put
\[
 G=G_S,\qquad \delta=m-e,\qquad B=m-1,\qquad
 \mathcal V=e+\frac\delta2-\frac32.
\]
The number $\mathcal V$ is a resource bound, not the width of $S$.
We prove the following assertion.

\begin{proposition}\label{prop:ordinary}
If $g_e-m\ge m-1$, then, over every field $k$,
\[
 \beta_i(G_S;k)\le\beta_i(G_{\widetilde S};k)
   =i\binom m{i+1}\quad(i\ge1).
\]
The zeroth Betti numbers on both sides are one.
\end{proposition}

\subsection{Coordinates, one cut, and one resource sequence}

The degree-$n$ monomial basis of $G$ consists
of $[v]$ with $\operatorname{ord}(v)=n$; multiplication survives precisely
when orders add.  Set
\[
 x=[m],\quad A=G/xG,\quad T=(0:_Gx),\quad
 S_0=k[Y_2,\ldots,Y_e].
\]
A basis element of $A$ is a \emph{head}: its coordinate has no surviving
incoming $x$-arrow.  Non-Ap\'ery heads span an $S_0$-ideal $N$ in $A$.
The quotient $Q=A/N$ has basis $\operatorname{Ap}(S,m)$ and is standard
graded.  For $q_s=\dim_kQ_s$, $a=q_2$, and
$\nu=\max\{s:q_s\ne0\}$, we have
\begin{equation}\label{eq:ordinary-apery}
 q_0=1,\quad q_1=e-1,\quad \sum_{s\ge2}q_s=\delta,
 \qquad A_2=Q_2,\quad \nu\le\delta+1.
\end{equation}
Indeed, Ap\'ery elements have maximal expressions without $m$, and every
prefix of such an expression realizes its order.  This proves standard
generation and the displayed counts; absence of gaps in the Hilbert
function gives the last inequality.

The monomial basis of $T$ consists of \emph{terminals} $\tau$ whose
outgoing $x$-arrow dies.  Their jumps
$d(\tau)=\operatorname{ord}(\tau+m)-\operatorname{ord}(\tau)$ satisfy
$2\le d(\tau)\le\nu$.  Translation $\tau\mapsto\tau+m$ bijects terminals
with non-Ap\'ery heads.  The appropriate module comparison uses the
opposite filtrations
\[
 D_sT=\langle\tau:d(\tau)\ge s\rangle,\qquad
 F_sN=\langle\tau+m:d(\tau)\le s\rangle.
\]
Both are $S_0$-filtrations, and their equal-jump strata are isomorphic:
\begin{equation}\label{eq:ordinary-strata}
 D_sT/D_{s+1}T\simeq F_sN/F_{s-1}N,
 \qquad\text{with degree shift }s.
\end{equation}
In fact, a surviving terminal action can only increase the jump, whereas
a surviving head action can only decrease it.  When the jumps agree, the
two actions survive together and have coefficient one.  To bound jumps,
write $Q=S_0/I_Q$ and $A=S_0/I_A$.  Since $Q_{\nu+1}=0$, $I_Q$, and hence
$N=I_Q/I_A$, is generated in degrees at most $\nu+1$.  A minimal head has
a positive-degree terminal, so its jump is at most $\nu$; head actions
propagate this bound.  A non-Ap\'ery head has degree at least three, which
also proves $A_2=Q_2$.

Let $p=\max\{u\ge0:\binom{u+1}{2}\le a\}$ and $r=\max(2,p)$.
We may replace $k$ by an infinite purely transcendental extension of the
same characteristic.  Green's degree-two restriction theorem
\cite{GreenRestriction} supplies independent linear forms
$E=(\lambda_1,\ldots,\lambda_r)$ such that $(A/EA)_{\ge2}=0$.
Indeed, if $a=\binom{p+1}{2}+b$, $0\le b\le p$, successive restrictions
bound the quadratic dimension after $u\le p$ cuts by
$\binom{p+1-u}{2}+\max(b-u,0)$, which vanishes at $u=p$.
If $p<2$, extend the cut to two forms.  Thus
\begin{equation}\label{eq:ordinary-cut-rank}
 2\le r\le e-1,\qquad
 r\ge3\Longrightarrow a_r:=\binom{r+1}{2}\le a\le\delta.
\end{equation}
Fix these forms once and put $P=k[\lambda_1,\ldots,\lambda_r]$.

For $2\le s\le\nu$, abbreviate $F_s=F_sN$ and define
\[
 U_s=A/F_{s-1},\quad X_s=(U_s)_{\ge s},\quad
 M_s=(U_s)_{\ge s+1},\quad K_s=F_s/F_{s-1}.
\]
Here $F_1=0$, $F_\nu=N$.  Jump-$s$ heads have degree at least $s+1$,
so these are actual exact sequences over the fixed ring $P$:
\begin{equation}\label{eq:ordinary-tower}
 0\longrightarrow K_s\longrightarrow M_s\longrightarrow X_{s+1}
 \longrightarrow0,\qquad X_{\nu+1}=0.
\end{equation}
The square-zero cut makes $U_s$ generated over $P$ in degrees zero and
one.  Its truncations $X_s,M_s$ are consequently generated only in
degrees $s,s+1$, respectively; their numbers of generators are $q_s$
and $\dim_k(X_s)_{s+1}$.
More explicitly, $(U_s)_n=E(U_s)_{n-1}$ for every $n\ge2$.
Iteration proves generation of each truncation in its bottom degree.
The degree-$s$ piece of $U_s$ agrees with $Q_s$, because every head
remaining outside $Q$ has jump at least $s$ and degree at least $s+1$.
When $\delta=0$, take $X_2=0$ in the formulas below.

\begin{lemma}\label{lem:ordinary-resources}
The finite sequence
$H_t=\sum_{s=2}^\nu\dim_k(X_s)_{s+t}$ satisfies, simultaneously,
\begin{equation}\label{eq:ordinary-resources}
 \begin{gathered}
 H_0=\delta,\qquad 0\le H_t\le B,\qquad
 S_H:=\sum_{t\ge0}H_t\le\delta\mathcal V,\\
 H_1=\delta-a+\ell\le B-a,\qquad H_1\le r\delta,
 \qquad \ell=\dim_kT_1.
 \end{gathered}
\end{equation}
For $\delta=0$ the tower and resource sequence are empty.
\end{lemma}

\begin{proof}
For a residue $\rho$ modulo $m$, let $b_\rho$ be the least sum of positive
offsets $g_i-m$ in that residue, and let $l_\rho$ be the shortest length
realizing $b_\rho$.  Write $\omega_\rho$ for the Ap\'ery representative,
$s_\rho=\operatorname{ord}(\omega_\rho)$, and
$L_\rho=(\omega_\rho-b_\rho)/m$.  Then
$1\le s_\rho\le L_\rho\le l_\rho$ for $\rho\ne0$.
Prefixes of a least-offset word are least-offset words and visit distinct
residues.  Along this path $L$ increases by at most one at each step,
and, after the first edge, cannot increase on arrival at an atom residue:
an increase there would decompose that atom. There are $e-1$ atom residues and $\delta$ other
nonzero residues.  Ordering nonzero residues by increasing $b_\rho$ gives
\[
 L_{\rho_j}-1\le\min(\delta,j-1),\qquad
 \sum_{\rho\ne0}(L_\rho-1)\le
 \sum_{j=1}^{m-1}\min(\delta,j-1)=\delta\mathcal V.
\]
Moreover, $v(q)=\operatorname{ord}(\omega_\rho+qm)$ eventually equals
$L_\rho+q$: the least-offset word gives this expression, and a longer
one would have smaller offset.  Telescoping therefore yields
\begin{equation}\label{eq:ordinary-terminal-budget}
 \sum_{\tau\equiv\rho\ (m)}(d(\tau)-1)=L_\rho-s_\rho.
\end{equation}
In $H$, an Ap\'ery coordinate contributes on $[0,s_\rho-2]$; a terminal
of degree $u$ and jump $d$ contributes its head on $[u,u+d-2]$.
These intervals are disjoint within each residue, since the next terminal
has degree at least $u+d$.  The zero residue contributes nothing.
Thus $H_t\le m-1$ and
$S_H=\sum_{\rho\ne0}(L_\rho-1)$.
At $t=0$ we count the $\delta$ Ap\'ery coordinates of degree at least two;
at $t=1$ we count $\delta-a$ higher Ap\'ery coordinates and $\ell$ linear
terminals.  Finally $\ell\le e-1$, and summing
$\dim_k(X_s)_{s+1}\le rq_s$ gives $H_1\le r\delta$.
\end{proof}

\subsection{From the tower to every relative row}

For a polynomial ring $\Lambda$, write
$b_i^\Lambda(Z)=\dim_k\operatorname{Tor}_i^\Lambda(Z,k)$ and
$c_j^\Lambda(Z)=\sum_{i=0}^j(-1)^{j-i}b_i^\Lambda(Z)$.
For a rank-zero module this is a syzygy rank and is nonnegative.

\begin{lemma}\label{lem:ordinary-harmonic}
If a finite-length $P$-module $Z$ is generated only in degree $b$ and
$h_t=\dim_kZ_{b+t}$, then
\[
 c_j^P(Z)\le\binom{r-1}{j}\sum_{t\ge0}\frac{j}{t+j}h_t
 \quad(1\le j<r),\qquad c_0^P(Z)=\mu_P(Z),\quad c_r^P(Z)=0.
\]
\end{lemma}

\begin{proof}
Shift to $b=0$ and take a homogeneous presentation
$P^q\twoheadrightarrow Z$. Choose a monomial order on its free module;
the monomials outside the initial module of the kernel give a basis of
$Z$, closed under divisors within each generator component. Label a
Koszul basis vector $m e_F$ by its augmented module monomial
$m\prod_{v\in F}x_v$. Reduction to standard monomials never increases
this label. Thus, in each internal degree, the actual differential is
block triangular in these labels, and its diagonal blocks are the
differential for the initial-module quotient. Choosing a maximal
nonsingular minor in each diagonal block and combining them gives a
block triangular minor of the actual matrix. Consequently its rank is
at least the sum of the diagonal ranks. The chain dimensions are
unchanged, so the identity for $c_j$ below reduces the desired upper
bound to standard monomials of monomial quotients.

In multidegree $\alpha$, $|\alpha|=t+j$, fix a variable $v$ in its support.
Every target row $(m,F)$ with $|F|=j$ and $v\notin F$ pairs with the
source column $(m/x_v,F\cup\{v\})$; divisors of standard monomials are
standard.  The selected minor is diagonal with entries $\pm1$.
Averaging these integer rank bounds with rational weights
$\alpha_v/(t+j)$ and summing gives
\[
 \operatorname{rank}\partial_{j+1}\big|_t\ge
 \binom{r-1}{j}\frac{t}{t+j}h_t.
\]
For $\lambda=\operatorname{length}Z$ and total boundary ranks $R_i$,
$c_j^P(Z)=\binom{r-1}{j}\lambda-R_{j+1}$, proving the claim.
The averaging uses rational numbers as numerical weights and hence
works in every characteristic.  The endpoint identities are the
generator count and the rank Euler identity.
\end{proof}

Set $R=k[x,E]$ and use $c_j(G)=c_j^R(G)$ from now on. The two-row $x$-Koszul
spectral sequence gives
$c_j(G)\le c_j^P(A)+c_{j-1}^P(T)$; its possible subsequent differential
subtracts a nonnegative rank.  For any exact sequence
$0\to K\to\mathcal M\to L\to0$, the Tor long exact sequence gives
\[
 c_j(\mathcal M)\le c_j(K)+c_j(L),\qquad
 c_j(K)\le c_j(\mathcal M)+c_{j+1}(L).
\]
For the first assertion, if $\rho_i$ denotes the rank of
$\operatorname{Tor}_i^P(A,k)\to\operatorname{Tor}_{i-2}^P(T,k)$, then
\[
 b_i^R(G)=b_i^P(A)+b_{i-1}^P(T)-\rho_i-\rho_{i+1},\qquad
 c_j(G)=c_j^P(A)+c_{j-1}^P(T)-\rho_{j+1}.
\]
For the second assertion, alternating the long exact sequence subtracts
either a connecting-map rank or its nonnegative complement in the
corresponding Tor group.  Thus all three comparisons preserve the
direction needed for an upper bound.
Apply these successively to $A$, the strata
\eqref{eq:ordinary-strata}, and \eqref{eq:ordinary-tower}.
Since $A/X_2\simeq P/P_+^2\oplus k(-1)^{e-r-1}$, we obtain
\begin{equation}\label{eq:ordinary-row}
 c_j(G)\le L_j+\sum_{t\ge0}W_j(t)H_t\quad(1\le j\le r),\qquad
 L_j=(e-1)\binom{r-1}{j}-\binom{r-1}{j+1},
\end{equation}
where
\[
 \begin{array}{ll}
 W_1(0)=r-1,&W_1(1)=(r+1)/2,\quad
 W_1(t)=(r-1)/(t+1)\quad(t\ge2),\\[2pt]
 W_j(0)=\binom{r-1}{j},&
 W_j(t)=\displaystyle\frac{j\binom{r-1}{j}}{t+j}
       +\frac{(j-1)\binom{r-1}{j-1}}{t+j-2}
       \quad(j\ge2,\ t\ge1).
 \end{array}
\]
Here the $X_s$ terms telescope, with last term $X_{\nu+1}=0$.
Lemma~\ref{lem:ordinary-harmonic} applies to $X_s$ and $M_s$;
for $j=1$ use the exact identity $\sum_s\mu_P(M_s)=H_1$.
The second denominator is shifted because $M_s$ starts one degree
later.  This also proves \eqref{eq:ordinary-row} at $j=r$.

\subsection{The target margin and the zeroth-homology correction}

Put $q=e-r$ and, with $e,r$ fixed, define
\[
 R_j(d)=(q+j)\binom{d+r-1}{j}
          +j\binom{d+r-1}{j+1},\qquad
 T_u(z)=1+\sum_{i\ge1}i\binom u{i+1}z^i.
\]
Pascal's identity gives
\[
 R_j(0)=L_j,\qquad
 R_j(d+1)-R_j(d)=(e+d)\binom{d+r-1}{j-1}.
\]
Consequently
\begin{equation}\label{eq:ordinary-margin}
 R_j(\delta)-L_j\ge
 \delta(\mathcal V+1)\binom{r-1}{j-1}.
\end{equation}
To explain why these are the required margins, extend $E$ to a basis
of the non-$x$ linear space, with $q-1$ remaining variables $Z$.
As a module over $R=k[x,E]$, $G$ has rank zero and $q$ generators.
Thus, writing $C(z)=\sum_{j=1}^rc_j(G)z^j$, its relative Betti polynomial
is $(1+z)(q+C(z))$.  Moreover
$G/(x,E)G\simeq k[Z]/(Z)^2$, whose Betti polynomial is $T_q$.
In the iterated Koszul complex, evaluate this zeroth homology row
exactly and bound the positive rows by their chain dimensions.  With
$\preceq$ denoting coefficientwise comparison, this gives
\begin{equation}\label{eq:ordinary-full-transfer}
 \mathcal B_G\preceq T_q+(1+z)^{q-1}
       \bigl((1+z)(q+C(z))-q\bigr)
       =T_{q+1}+(1+z)^qC(z).
\end{equation}
Here the relative degree-zero homology has dimension $q$, but the
remaining variables act on it: it is the square-zero quotient
$k[Z]/(Z)^2$, not a trivial $q$-dimensional coefficient space. Its
contribution must therefore be $T_q$, rather than
$q(1+z)^{q-1}$. Subtracting $q$ inside the parentheses removes the
degree-zero term before the positive relative rows are bounded by
chain dimensions. This prevents double counting and gives the exact
zeroth row $1$. No spectral-sequence degeneration is used.
Summing $T_{u+1}-T_u=uz(1+z)^{u-1}$ yields
\begin{equation}\label{eq:ordinary-target-difference}
 T_m-T_{q+1}=(1+z)^q
       \sum_{j=1}^{\delta+r-1}R_j(\delta)z^j.
\end{equation}
It therefore suffices to prove $c_j(G)\le R_j(\delta)$ for $1\le j\le r$;
the additional target coefficients are nonnegative.
Indeed the coefficient before factoring out $(1+z)^q$ is obtained from
$\sum_{u=0}^{\delta+r-2}(q+1+u)\binom{u}{j-1}$.
The hockey-stick identity and
$u\binom{u}{j-1}=(j-1)\binom{u}{j-1}+j\binom{u}{j}$ evaluate it as
$R_j(\delta)$.  This calculation includes $\delta=0$: out-of-range
binomial coefficients vanish, and the relative top margin is zero.

\subsection{The numerical inequalities}

For $r=2$, \eqref{eq:ordinary-row} and
$H_1\le\min(S_H-\delta,2\delta)$ give
\[
 c_1-L_1\le\min\{\tfrac32S_H-\tfrac12\delta,
                         \tfrac52\delta+\tfrac12S_H\}
 \le\delta(\mathcal V+1).
\]
Use the first estimate for $\mathcal V\le3$ and the second for
$\mathcal V\ge3$.  The top row satisfies $c_2\le S_H-\delta$.
For $r=3$, $e\ge4$, $\delta\ge6$, and $\mathcal V\ge11/2$.
The first row is at most
$L_1+\frac43\delta+\frac43H_1+\frac23S_H$.
Using $H_1\le B-6$, three times its margin is bounded below by
\[
 \delta(\mathcal V-1)-4(B-6)
 =zs+\tfrac12z^2+\tfrac72z+2s+15>0,
 \quad z=\delta-6,\ s=e-4.
\]
The second row is at most
$L_2+\frac32S_H-\frac12\delta+\frac76H_1
 \le L_2+\frac32S_H+3\delta
 \le L_2+2\delta(\mathcal V+1)$;
the top row is again at most $S_H-\delta$.
Thus \eqref{eq:ordinary-margin} proves all low-rank comparisons.

For $r\ge4$ define the single scalar
\[
 F_r=(r-1)\delta+\frac{r+1}{2}H_1
              +\sum_{t\ge2}\frac{r-1}{t+1}H_t.
\]
Dividing $W_j(t)$ by $\binom{r-1}{j-1}$ bounds it by the corresponding
coefficient of $F_r$.  At $t\ge2$ the normalized weight is
$(r-j)/(t+j)+(j-1)/(t+j-2)$.
For $j\ge3$ both denominators are at least $t+1$; for $j=2$ its difference
from $(r-1)/(t+1)$ is
$[2-(r-3)t]/[t(t+1)(t+2)]\le0$.
The cases $t=0,1$ follow directly from the binomial identities.
These estimates depend only on the displayed resource inequalities.
Hence all rows, including the top, follow once
\begin{equation}\label{eq:ordinary-scalar}
 F_r\le\delta(\mathcal V+1).
\end{equation}

Here is a common cutoff argument for this last inequality.
Set $a_*=a_r$ and choose an integer $n\ge r$; write
$\alpha=(r-1)/(n+1)$, $\beta=1-\alpha$,
$c=(r+1)/2-\alpha$, and $\mathsf h_n=\sum_{j=1}^n1/j$.
Subtracting $\alpha$ from every weight and using
\eqref{eq:ordinary-resources} gives
\begin{equation}\label{eq:ordinary-cutoff}
 \begin{split}
 F_r\le U(e,\delta):={}&\alpha\delta\mathcal V
 +(r-1-\alpha)\delta+c(B-a_*)\\
 &+B\sum_{j=2}^{n-1}\left(\frac{r-1}{j+1}-\alpha\right).
 \end{split}
\end{equation}
The omitted residual weights for $j\ge n$ are nonpositive; no support
bound is assumed.  Define $\Delta(e,d)=d(e+d/2-1/2)-U(e,d)$.
For any base $e_0$, put $g=\Delta(e_0,a_*)$,
$A_*=\Delta(e_0+1,a_*)-g$, and $b=\Delta(e_0,a_*+1)-g$.
Summing the finite harmonic expression gives
\begin{equation}\label{eq:ordinary-gap-shift}
 \begin{gathered}
 A_*=a_*-1-(r-1)\left(\mathsf h_n-2+\frac{a_*+2}{n+1}\right),\\
 \Delta(e_0+s,a_*+t)
 =g+sA_*+tb+\beta\left(\frac{t(t-1)}2+st\right).
 \end{gathered}
\end{equation}
This is an identity in rational $s,t$.

For $4\le r\le8$ use $e_0=r+1$ and the following exact values:
\[
 \begin{array}{c|r|r|r|r|r|r}
 r&a_*&n&L&Lg&LA_*&Lb\\\hline
 4&10&7&280&2443&762&1007\\
 5&15&10&6930&39380&28429&32209\\
 6&21&13&72072&-7479&424147&439591\\
 7&28&17&2042040&-12196485&17076937&17076937\\
 8&36&21&739024&-5070732&8418651&8250691
 \end{array}
\]
All $A_*,b,\beta$ are positive.  The first two seeds are positive;
in the last three rows $g+A_*,g+b>0$.  Thus
\eqref{eq:ordinary-gap-shift} is positive for nonnegative integer $s,t$,
except possibly $s=t=0$ in the last three rows.
Those corners cannot occur.  Indeed, if $r=e-1\ge3$, then
$a=\binom e2$, so all quadratic products of the non-$m$ atoms give
distinct nonzero degree-two Ap\'ery coordinates.  The residues of
$\{0,g_2,\ldots,g_e\}$ consequently form a Sidon set modulo $m$.
Its $e(e-1)$ ordered nonzero differences are distinct, giving
$m\ge e(e-1)+1$ and $\delta\ge(e-1)^2>a_*$.
This proves \eqref{eq:ordinary-scalar} in these ranks.

For completeness, the unbounded range has a uniform argument.
Let $r\ge9$, $a_*=r(r+1)/2$, and $D=\lfloor a_*/2\rfloor+1$;
use $n=D$ in \eqref{eq:ordinary-cutoff}.  Then
$a_*+1\le2D\le a_*+2$ and $D\ge r$.
For $\mathsf b(r)=r/4+1+9/(2r)$,
\[
 \mathsf h_{D_{r+1}}-\mathsf h_{D_r}
 \le\frac{r+2}{a_r+3}<\mathsf b(r+1)-\mathsf b(r)\quad(r\ge10).
\]
The strict inequality follows on putting $x=r-10$ from the positive
numerator $x^4+34x^3+385x^2+1472x+112$ over
$4r(r+1)(r^2+r+6)$.
Together with
$\mathsf b(10)-\mathsf h_{28}=1833472237/80313433200>0$, this proves
$\mathsf h_D<\mathsf b(r)$ for every $r\ge10$.
Define
\[
 T_r=\frac{a_*(a_*+22-r)}{2(r-1)(a_*+9)}
       +\frac{2r-3}{r-1}
       -\frac{a_*(a_*+21)+36}{2(D+1)(a_*+9)}.
\]
Direct summation in \eqref{eq:ordinary-cutoff} gives
$g_0:=\Delta(10,a_*)=(r-1)(a_*+9)(T_r-\mathsf h_D)$.
Replace the last denominator of $T_r$ by $(a_*+3)(a_*+9)$ to obtain
$T_-\le T_r$.  Its difference from $\mathsf b(r)$ is
\[
 T_--\mathsf b(r)=
 \frac{r^3[(r-10)(2r+15)+2]+r[111(r-10)+210]+972}
 {2r(r-1)(r^2+r+6)(r^2+r+18)}>0.
\]
Thus $g_0>0$ for $r\ge10$; for $r=9$, direct substitution gives
$g_0=5845533/7436429>0$.

It remains to extend this seed in both parameters.
With $A_*=\Delta(11,a_*)-g_0$, $b_0=\Delta(10,a_*+1)-g_0$, put
\[
 C_*=\frac{a_*[-a_*(D-r+2)-D(r-4)+3]}{2(D+1)}.
\]
The identities $g_0=(a_*+9)A_*+C_*$ and
$b_0=A_*+11-r-8\alpha$ give $A_*>-C_* /(a_*+9)>0$.
For $F(D)=-C_* /(a_*+9)+11-r-8\alpha$, set
\[
 F_{\min}=\frac{a_*}2-\frac{3r}2+\frac{11}2
       +\frac{3(5r+31)}{2(a_*+9)}-\frac{17(r-1)}{a_*+3}.
\]
The floor comparison is the identity
\[
 F(D)-F_{\min}=
 \frac{(2D-a_*-1)(r-1)(a_*^2+17a_*+144)}
 {2(a_*+9)(D+1)(a_*+3)}\ge0.
\]
Since $a_*\ge4r$, $F_{\min}>11/2-17/4=5/4$;
hence $b_0>5/4$.  Finally $0<\beta<1$, and
\eqref{eq:ordinary-gap-shift}, recentered at $e_0=10$, reads
\[
 \Delta(10+s,a_*+t)=g_0+sA_*+t(b_0-\beta/2)
                       +\beta(t^2/2+st)>0
 \qquad(s,t\ge0).
\]
This proves the scalar inequality for all rational $e\ge10$,
$\delta\ge a_*$, and therefore for every actual cut with $r\ge9$.

\begin{proof}[Proof of Proposition~\ref{prop:ordinary}]
The low-rank estimates and the scalar inequality, combined with
\eqref{eq:ordinary-row} and \eqref{eq:ordinary-margin}, give
$c_j(G)\le R_j(\delta)$ for every $1\le j\le r$.
Equations~\eqref{eq:ordinary-full-transfer}--\eqref{eq:ordinary-target-difference}
yield $\beta_i(G;k)\le i\binom m{i+1}$.
The width assumption identifies the minimal generators of the interval
completion as $m,m+1,\ldots,2m-1$; its Betti numbers are those in
\eqref{eq:ordinary-capacity}, by \cite{APFormula,APCone}.
Finally, scalar extension preserves the original monomial coordinates
and the ranks in a minimal graded free resolution.  Thus the bound
proved after the purely transcendental extension descends to the
original field.  The algebraic and completed tangent cones have the
same graded coordinates by Proposition~\ref{prop:bridge}; the conclusion
therefore concerns the stated completed semigroup rings and their own
minimal presentations.
\end{proof}

\section{The endpoint mechanism}\label{sec:endpoint}

Let $k$ be the original coefficient field, and suppose
\[
 q\ge2,\qquad 0<a_1<\cdots<a_q<u<m,\qquad d\ge1,
 \qquad \gcd(a_1,\ldots,a_q,u)=\gcd(m,d)=1.
\]
Set
\[
 \Gamma=\langle m,m+da_1,\ldots,m+da_q,m+du\rangle,
 \qquad G=\operatorname{gr}_{\mathfrak m}k[[t^\Gamma]],
\]
and write $S=k[x,y_1,\ldots,y_q,v]$, $E=k[x,v]$ and
$T=k[y_1,\ldots,y_q]$, with every variable of degree one.
We use $\beta_i(G)$ over $S$ and $b_i(A)$ over $T$.
Binomial coefficients and Betti numbers with negative subscripts are zero.
The physical width is $V=du$; $u$ is the normalized width.
In particular, normalization retains $m$ and $d$ and describes the same
semigroup and tangent cone. The endpoint construction needs $u<m$;
the narrow-width application also has $V\le m-2$.

\paragraph{The core and its fibers.}
Consider the standard graded algebra
\[
 C=k[s,st^{a_1},\ldots,st^{a_q},st^u],
 \qquad x=s,\quad y_j=st^{a_j},\quad v=st^u.
\]
The displayed generators of $\Gamma$ are minimal. Indeed, a decomposition
of the generator with offset $a\in\{0,a_1,\ldots,a_q,u\}$ into
$l\ge2$ generators, with offset sum $B$, would give
$(l-1)m=d(a-B)$. Coprimality forces $l-1=jd$ and $a-B=jm$
for some $j\ge1$, contradicting $a\le u<m$.
If $\nu(h)$ denotes the maximum length of a generator expression for $h$,
there is a graded surjection
\begin{equation}\label{eq:endpoint-core-map}
 C\longrightarrow G,\qquad
 s^nt^B\longmapsto
 \begin{cases}(t^{nm+dB})^*,&\nu(nm+dB)=n,\\
 0,&\nu(nm+dB)>n.
 \end{cases}
\end{equation}
Maximum-length expressions prove surjectivity and compatibility with products.
Distinct offsets in a fixed degree give distinct numerical exponents, so
its kernel is spanned by the core monomials that die.

For $0\le r<u$, put $\epsilon_r=0$ if $r=0$ and $\epsilon_r=1$ otherwise.
Write the degree and offset of a core monomial as
\[
 B=r+uK,\qquad n=L+K+\epsilon_r.
\]
Since $B\le nu$, both $L$ and $K$ are nonnegative.
Multiplication by $x$ or $v$ increases respectively $L$ or $K$.
Let $M_r\subseteq\N^2$ be the resulting upward closed set, also viewed
as a monomial ideal of $E$. Thus, as $E$-modules,
\begin{equation}\label{eq:endpoint-fibers}
 C\simeq\bigoplus_{r=0}^{u-1}M_r,
\end{equation}
with the degree of $(L,K)$ in the $r$th summand equal to
$L+K+\epsilon_r$.

These ideals have a finite, intrinsic description. A word of at least $u$
internal letters has two prefix sums equal modulo $u$. The intervening
nonempty word has length $l$ and offset $cu$, where $1\le c<l$;
replace it by $x^{l-c}v^c$. Repetition leaves at most $u-1$ internal
letters. Consequently every minimal point of $M_r$ is represented by a
pure internal word of length at most $u-1$. All residues occur, because
the internal offsets generate $\mathbb Z/u\mathbb Z$ and inverses in this
finite group can be expressed by nonnegative repetitions. Moreover,
$M_0=E$, and every other minimal point satisfies $L+K\le u-2$.
In particular $C$ is finite over $E$, and
\[
 H_x:=\max\{L:(L,K)\in\operatorname{Min}M_r\text{ for some }r\}
 \le u-2.
\]

\paragraph{All competitors and finite generation.}
Equality of numerical exponents shows that a longer expression for a
core monomial $(n,B)$ is exactly a core point
\[
 (n+jd,B-jm),\qquad j\ge1.
\]
Fix the longer-side residue $r$ and a minimal point $(L,K)$ of $M_r$.
Write $r+jm=s+cu$, with $0\le s<u$. On the shorter side this
competitor forbids the quadrant starting at
\begin{equation}\label{eq:endpoint-competitor}
 (A,B')=
 \bigl(\max\{0,L-jd-c+\epsilon_r-\epsilon_s\},K+c\bigr).
\end{equation}
The maximum accounts exactly for increasing the longer-side $x$ exponent
until the shorter-side first coordinate is nonnegative.
Let $T_s\subset E$ be generated by these starts for every $j\ge1$,
every $r$, and every minimal point of $M_r$. Then
\begin{equation}\label{eq:endpoint-kernel}
 N_s=M_s\cap T_s,\qquad
 G\simeq\bigoplus_{s=0}^{u-1}M_s/(M_s\cap T_s)
\end{equation}
as $E$-modules. The intersection imposes membership in the core;
the competitor ideal alone need not be the kernel.

Although all positive $j$ enter this description, the generating process
is finite. Indeed,
\[
 c-\epsilon_r+\epsilon_s
 =\left\lceil\frac{r+jm}{u}\right\rceil
  -\left\lceil\frac r u\right\rceil
 \ge\left\lfloor\frac{jm}{u}\right\rfloor\ge j,
 \qquad
 A\le\max\{H_x-2j,0\}.
\]
Thus all starts have first coordinate at most $H=\max\{H_x-2,0\}$.
For each first coordinate that occurs, retain only the smallest second
coordinate; these finitely many starts generate $T_s$.
Choosing $j$ with $2j\ge H_x$ and $r\equiv s-jm\pmod u$ also gives a
start with first coordinate zero. Hence every $T_s$ contains a pure
$v$ power. This proves the finite description directly from all
competitors, without a coprimality assumption on $m$ and $u$.

\paragraph{Length and socle resources.}
Define
\[
 \begin{gathered}
 Q_C=C/(x,v)C,\qquad Q_G=G/(x,v)G,\\
 P_C=\dim_k Q_C,\quad P_G=\dim_k Q_G,\quad
 R=\dim_k(0:_G(x,v)).
 \end{gathered}
\]
The minimal fiber points form a basis of $Q_C$, and their surviving
subset forms a basis of $Q_G$. As a $T$-algebra $Q_C$ is standard graded
Artinian, with Hilbert function
$h_n=\dim_k(Q_C)_n$ satisfying $h_0=1$, $h_1=q$ and $h_n=0$ for
$n\ge u$. It need not be a monomial quotient of $T$.
Let $\ell_r$ be the shortest internal word length reaching $r$ modulo $u$,
with $\ell_0=0$, and put
\begin{equation}\label{eq:endpoint-resources}
 W=1+\sum_r\ell_r,\qquad
 \widehat R=\min\{W-u,u\max\{H_x-2,0\}\}.
 \qquad P_C\le W,\quad R\le\widehat R.
\end{equation}

Here is the counting argument. A shortest word determines a minimal
point $(L_r,K_r)$. In the ordered staircase of minimal points, at most
$L_r$ points lie to its left and at most $K_r$ to its right. For
$r\ne0$ this gives at most $L_r+K_r+1=\ell_r$ points; residue zero
contributes one. This proves $P_C\le W$ and also $W\ge u$.

Order the minimal generators of $T_s$ as
\[
 (A_0,B_0),\ldots,(A_z,B_z),\qquad
 0=A_0<\cdots<A_z,\quad B_0>\cdots>B_z=\kappa_s.
\]
The socle of $E/T_s$ has the $z=z_s$ corner monomials
$(A_{j+1}-1,B_j-1)$ as a basis. This assertion only concerns its socle;
$E/T_s$ need not have finite length. Counting generators on either side
of a minimal point, and then passing to any dominating point, gives
\[
 z_s\le A+B-\kappa_s\quad((A,B)\in T_s),\qquad
 \kappa_s\ge\left\lceil\frac{m-s}{u}\right\rceil.
\]
The natural injection $M_s/(M_s\cap T_s)\hookrightarrow E/T_s$ induces
a socle injection, so $R\le\sum_s z_s$.
Apply the $j=1$ competitor to each shortest point, and set
$s_r=(r+m)\bmod u$, $c_r=(r+m-s_r)/u$.
Its start has $A\le L_r$ and $B=K_r+c_r$.
Translation permutes the residues, and
\[
 \sum_r c_r=m=\sum_s\left\lceil\frac{m-s}{u}\right\rceil.
\]
Summing the point bounds therefore yields
$\sum_s z_s\le\sum_r(L_r+K_r)=W-u$.
Independently, the first coordinates of minimal generators lie in
$[0,H]$, whence $z_s\le H$. Together these prove
\eqref{eq:endpoint-resources}. Every resource refers to the same
orientation of the core; replacing $H_x$ by a reflected height would
require a separate comparison with this $R$.

\paragraph{The compensated Tor bound.}
The endpoint Koszul homology is finite-dimensional and has dimensions
\begin{equation}\label{eq:endpoint-homology}
 \bigl(\dim_k H_0(x,v;G),\dim_k H_1(x,v;G),
       \dim_k H_2(x,v;G)\bigr)=(P_G,P_G+R,R).
\end{equation}
For completeness, $G/xG$ and $(0:_Gx)$ have finite dimension by the
Ap\'ery description of the multiplicity action. View
$K(x,v;G)$ as the cone of multiplication by $v$ on $K(x;G)$.
The cone long exact sequence, applied to these two finite homology
groups, gives the middle equality in \eqref{eq:endpoint-homology} by
rank--nullity. The zeroth and second groups are respectively $Q_G$ and
$(0:_G(x,v))$. No finite-dimensionality assertion about $G$ is involved.

Let $K$ be the kernel of the actual $T$-linear surjection $Q_C\to Q_G$.
Then $\dim_k K=P_C-P_G\ge0$. The Tor long exact sequence and the
dimensions of the internal Koszul chain groups give
\[
 b_i(Q_G)\le b_i(Q_C)+(P_C-P_G)\binom q{i-1}.
\]
In the bounded double Koszul complex on the actual $G$, take endpoint
homology first. Its zeroth row retains the genuine internal Tor of
$Q_G$; bounding the other two rows by their chain dimensions gives
\[
 \beta_i(G)\le b_i(Q_G)+(P_G+R)\binom q{i-1}
                              +R\binom q{i-2}.
\]
This is a dimension inequality for a convergent filtration, and requires
no degeneration assertion. Substitution cancels $P_G$ against the
quotient-kernel length. Pascal's identity yields
\begin{equation}\label{eq:endpoint-bound}
 \boxed{\displaystyle
 \beta_i(G)\le b_i(Q_C)+P_C\binom q{i-1}
                              +R\binom{q+1}{i-1}.}
\end{equation}
In particular $\beta_0(G)=1$, $\beta_{q+2}(G)\le R$, and
$\beta_i(G)=0$ for $i>q+2$.

\paragraph{Artinian estimates and the exported rows.}
For a standard graded Artinian $A=T/I$ with $I_1=0$, write
$P=\dim_k A$ and $h_n=\dim_k A_n$.
The cyclic, degree-zero generated case of
Lemma~\ref{lem:ordinary-harmonic} gives
\begin{align}
 b_1(A)&\le q+\sum_{n\ge1}\frac{(q-1)h_n}{n+1}
 \le \frac{q-1}{3}P+\frac{(q+1)(q+2)}6,
 \label{eq:endpoint-artin-first}\\
 b_1(A)&\le\frac{q-1}{4}P+\frac{(q+1)^2}{4}
                         +\frac{q-1}{12}h_2.\nonumber
\end{align}
To see the specialization, the first partial Euler sum is
$c_1(A)=b_1(A)-b_0(A)=b_1(A)-1$.
The lemma bounds it by
$(q-1)\sum_{n\ge0}h_n/(n+1)$, and $h_0=1$ supplies the constant
$q$ in \eqref{eq:endpoint-artin-first}. Its hypotheses hold because
the class of $1$ generates $A$ in degree zero and the homogeneous
pieces have finite support.
The second and third bounds separate degree one, or degrees one and two,
and bound the remaining reciprocal weights by $1/3$ or $1/4$.
Subtracting $P$ from the harmonic bound gives
\[
 b_1(A)-P\le q-1+\sum_{n\ge1}\frac{q-n-2}{n+1}h_n.
\]
For $q=5$, discard the nonpositive terms with $n\ge3$ and use
$h_2\le15$ to obtain $b_1(A)-P\le14$.
All these weights and signed differences are interpreted in $\mathbb Q$.

There is also a uniform estimate for every positive row:
\[
 b_i(A)+\binom{q+1}{i+1}\le P\binom qi\qquad(i\ge1).
\]
Indeed, in the Koszul differential $d_i$, restrict the source to
coefficient degree zero. For each wedge $J$, select the target coordinate
$(y_{\min J},J\setminus\{\min J\})$. These coordinates give a diagonal
minor with entries $\pm1$, since the variable images in $A_1$ are
independent. Hence $\operatorname{rank}d_i\ge\binom qi$; apply the
same argument to $d_{i+1}$ and subtract both ranks from the chain
dimension $P\binom qi$.
The minors remain invertible in every characteristic, including two.

Applying these estimates to $Q_C$ in \eqref{eq:endpoint-bound} gives
\begin{equation}\label{eq:endpoint-rows}
 \boxed{\begin{aligned}
 \beta_1(G)&\le P_C+
  \left\lfloor q+\sum_{n\ge1}\frac{(q-1)h_n}{n+1}\right\rfloor
  +\widehat R,\\
 \beta_i(G)+\binom{q+1}{i+1}
 &\le P_C\binom{q+1}i+\widehat R\binom{q+1}{i-1}
 \qquad(i\ge1).
 \end{aligned}}
\end{equation}
The floor is valid because the Betti number is integral.
These estimates concern the original completed tangent cone over $k$;
no extension of the coefficient field is needed here.

Finally suppose $q=2$, with internal offsets $a<b$ and variables $y,z$.
In fixed degree $n$, the words $y^{n-j}z^j$ have distinct offsets
$na+j(b-a)$, so $Q_C$ is a monomial quotient of $k[y,z]$.
Put $\sigma_C=\dim_k(0:_{Q_C}(y,z))$.
The two-variable Koszul complex and its Euler identity give
$(b_0(Q_C),b_1(Q_C),b_2(Q_C))=(1,\sigma_C+1,\sigma_C)$.
Thus \eqref{eq:endpoint-bound} specializes to
\begin{equation}\label{eq:endpoint-two-rows}
 \boxed{\begin{aligned}
 \beta_1(G)&\le P_C+\sigma_C+1+R,&
 \beta_2(G)&\le2P_C+\sigma_C+3R,\\
 \beta_3(G)&\le P_C+3R,&\beta_4(G)&\le R.
 \end{aligned}}
\end{equation}
Each $R$ may be replaced by $\widehat R$. The internal socle
$\sigma_C$ and the endpoint socle $R$ belong to different quotients
and different variable actions.

The core, competitor, resource and transfer arguments also apply with
one internal offset. For $q=1$, coprimality says $\gcd(a_1,u)=1$.
The relation $y^u=x^{u-a_1}v^{a_1}$ in $C$ shows that $Q_C$ is
spanned by the classes of $1,y,\ldots,y^{u-1}$, so $P_C\le u$.
The one-variable Artinian presentation is
$Q_C\simeq k[y]/(y^{P_C})$, whose Betti numbers are $(1,1)$.
Formula \eqref{eq:endpoint-bound} therefore gives
\[
 (\beta_1(G),\beta_2(G),\beta_3(G))
 \le (P_C+R+1,P_C+2R,R)
\]
coordinatewise. These rows use the same actual endpoint socle $R$;
any sharper bound for that socle can be substituted before comparing
with the interval target.

\section{The narrow-width comparison}\label{sec:width}

\begin{proposition}\label{prop:narrow}
For every numerical semigroup $S$ of multiplicity $m$, embedding
dimension $e\ge4$, and width at most $m-2$, every field $k$, and every $i\ge0$,
\[
 \beta_i(G_S;k)\le\beta_i(G_{\widetilde S};k).
\]
Both Betti numbers are taken over the polynomial ring on the respective
semigroup's minimal generators.
\end{proposition}

We prove the proposition by covering the possible normalized widths.
Let $S$ have multiplicity $m$, embedding dimension $e\ge4$, and width
$V\le m-2$.  Write its minimal generators as
\[
 S=\langle m,m+da_1,\ldots,m+da_q,m+du\rangle,
 \qquad q=e-2,\quad 0<a_1<\cdots<a_q<u,
\]
where $V=du$, $d\ge1$, $m\ge du+2$, and
$\gcd(a_1,\ldots,a_q,u)=\gcd(m,d)=1$.
This is a change of coordinates for the given generators: their
numerical values and the coefficient field stay fixed.  The comparison
target is still generated by all integers from $m$ through $m+du$.
In particular, any estimate proved at normalized width $u$ must be
transferred to that physical interval.
Use the core, height, cyclic distances, and invariants of
Section~\ref{sec:endpoint}, with $\Gamma=S$.
For brevity put $G=G_S$, $\beta_i=\beta_i(G;k)$,
$Q=Q_C$, $P=P_C$, and $h_n=\dim_k Q_n$.

The endpoint estimates of Section~\ref{sec:endpoint} give $P\le W$,
$R\le\widehat R$, and, writing $b_i=\beta_i(Q;k)$,
\begin{align}
 \beta_i&\le b_i+P\binom q{i-1}+R\binom{q+1}{i-1},\label{eq:width-transfer}\\
 \beta_i+\binom{q+1}{i+1}
 &\le P\binom{q+1}i+\widehat R\binom{q+1}{i-1}.
 \label{eq:width-rows}
\end{align}
The use of the core dimension $P$ in these inequalities includes the
kernel compensation from the endpoint argument; the actual endpoint
quotient can have a different dimension.  Rational coefficients in the
estimates below are weights on dimensions, so the calculations apply
over every characteristic.
The interval completion has, in row $i$, at least
\[
 F_i^-(u)=i\binom u{i+1}
       +\min\left\{\binom u{i-1},\binom ui\right\},
 \qquad 1\le i\le u.
\]
Indeed its physical width is $V\ge u$, and this lower bound increases
with width.  We therefore compare with $F_i^-(u)$, except at the full
arithmetic progression, where a direct comparison is simpler.
Binomial coefficients with an index outside their range are zero.

\subsection{The two smallest widths}

If $u=q+1$, every offset occurs.  The arithmetic progression formula
\cite{APFormula,APCone} bounds every positive row by
$i\binom{u+1}{i+1}$: its correction to $i\binom u{i+1}$ is at most
$i\binom ui$.  For $d\ge2$, $V\ge2u\ge u+1$ supplies this bound.
For $d=1$, $S$ already is its interval completion.

Suppose next that $u=q+2$ and $q\ge3$, and let $h$ be the missing
offset.  The case $q=2$ will be included in the four-generator argument.
For
$B=\{0,\ldots,u\}\setminus\{h\}$, a basis of $Q_n$ is indexed by
\[
 nB\setminus\bigl((n-1)B\cup((n-1)B+u)\bigr).
\]
If $2\le h\le u-2$, then $B+B=[0,2u]$: the missing sums $h,h+u$
are $1+(h-1)$ and $(u-1)+(h+1)$.  Thus $Q$ has Hilbert function
$(1,q,2)$, with quadratic basis at offsets $h,h+u$.  If $h=1$ or
$u-1$, the Hilbert function is $(1,q,1)$.  In every case these basis
elements have height at most one, so $R=0$.

For an interior hole, split the linear variables into blocks of sizes
$r=h-1$ and $s=u-h-1$, so $r+s=q$.  Multiplication pairs each variable
with its unique partner in the same block, producing the corresponding
quadratic basis element.  In the Koszul differential, every nonzero
quadratic row admits a distinct column with its only nonzero entry
$\pm1$ in that row.  Counting the rows whose wedges contain a whole
block, combining this minor with the degree-zero minor, and applying
\eqref{eq:width-transfer} gives
\[
 \beta_i+\delta\binom{q+1}i\le i\binom{q+2}{i+1}+T_i,
 \qquad
 T_i=\binom{s+1}{i-r}+\binom{r+1}{i-s}.
\]
Here $\delta\ge0$ is any integer satisfying $P+\delta\le q+3$;
the same bound allows an empty block at an end hole.
For the following binomial estimates assume an interior hole, so
$r,s\ge1$. End holes are treated separately below.
For $0\le c\le r$, adjoining $r-c$ fixed elements and then enlarging
the ambient set gives
\[
 \binom{s+1}{i-r}\le\binom{q+1}{i-c}.
\]
Apply this with $c=0$ in the first summand of $T_i$, and with $c=1$
in the second.  Pascal gives $T_i\le\binom{q+2}i$.
For $i\ge2$, use $c=2$ in a block of size at least two and $c=1$
in the other block.  Since $r+s\ge3$, this proves
$T_i\le\binom{q+2}{i-1}$; the case $i=1$ is immediate.
Taking $\delta=0$ proves $\beta_i\le F_i^-(u)$ for an interior hole.
For an end hole, without using these binomial estimates, take
$\delta=1$, since $P=q+2$.
The empty-block correction equals $\binom{q+1}i$
plus the two top-row indicators, giving
\[
 \beta_i\le i\binom{q+2}{i+1}
                   +\mathbf1_{i=q}+\mathbf1_{i=q+1}.
\]
Both exceptional rows have target correction at least one.  The minors
use units $\pm1$, so the argument includes characteristic two.

\subsection{Dense widths}

Assume $q\ge4$ and $u=q+1+t$, where $2\le t\le q$.
The prefixes of lengths $2,\ldots,n$ of a surviving word must have
distinct missing residues.  An existing residue would replace a prefix
by a word containing an endpoint; repeated residues would give a
nonempty zero-sum subword and the same contradiction.  Hence
$n-1\le t$ and $H_x\le t$.
The set $A=\{0,a_1,\ldots,a_q\}$ has more than half the residues, so
$A+A$ is the whole group.  Its distance layers consequently give
\begin{equation}\label{eq:width-dense-resources}
 W=q+1+2t,\quad P\le q+1+2t,\quad R\le t,
 \qquad t=2\Longrightarrow R=0.
\end{equation}

Choose one lex quotient with the entire Hilbert function of $Q$.
For this quotient let $S_j(n)$ count degree-$n$ standard monomials in
the first $j$ variables.  Counting minimal generators by their largest
variable gives
$g_j(n)=S_j(n-1)-S_j(n)+S_{j-1}(n)$.
The lex Betti bound, including Pardue's arbitrary-characteristic version
\cite[Theorem~31]{Pardue}, and the stable ideal formula \cite{BHP,EK}, followed
by telescoping in degree, imply
\begin{equation}\label{eq:width-lex}
 b_i\le i\binom{q+1}{i+1}+K\binom{q-1}{i-1},
 \qquad K=\sum_{n\ge2}\sum_{j=1}^qS_{j-1}(n).
\end{equation}
If $\binom{r+1}2\ge2t$ and $r\le q$, all higher standard monomials
are supported on the last $r$ variables.  Indeed each $h_n\le2t$,
and the degree-$n$ monomials in those variables already number at least
$\binom{r+1}2$ and come first among the lex standard monomials.
Count degree-$n+1$ monomials
whose direct divisors are standard, weighting a parent--child edge by
the exponent of the removed variable.  Each child has total weight
$n+1$ and each parent at most $n+r$.
Removing the largest variable from a child gives the prefix-standard
parent count $\sum_jS_j(n)$; its excess over $h_n$ is exactly the
degree-$n$ summand of $K$.  Consequently
\[
 K\le\sum_{n\ge2}\frac{r-1}{n+1}h_n
       \le\frac{2t(r-1)}3.
\]
These weights are rational numbers used to count finite sets.

Put $A_i=\binom{q-1}{i-1}$, $C_i=\binom q{i-1}$,
$B_i=\binom{q+1}{i-1}$.  The first two terms of Vandermonde give
\[
 i\left[\binom{q+t+1}{i+1}-\binom{q+1}{i+1}\right]
 \ge t(q+1)C_i+i\binom t2B_i.
\]
For $t=2$, choose $r=3$, giving $K\le2$; the remaining difference is
at least $(q-3)C_i+iB_i-2A_i\ge0$.
For $t\ge3$, choose $r=t$.  After using $B_i\ge C_i\ge A_i$, the
difference is at least $cC_i$, where
\[
 6c=12(q-4)+5(t-3)^2+(t-3)+6(t-3)(q-t)\ge0.
\]
Equations \eqref{eq:width-transfer}--\eqref{eq:width-lex} therefore
give $\beta_i\le i\binom u{i+1}$ in every positive row.

\subsection{Sparse widths and propagation from the second row}

We first record why one numerical comparison controls all higher rows.
For $e\le u$ define
\[
 U_i(B)=B\binom ei-u\binom{e-1}{i-1}-\binom{e-1}{i+1},
 \quad
 \mathcal T_i(u)=
 \frac{i\binom u{i+1}+u\binom{e-1}{i-1}+\binom{e-1}{i+1}}{\binom ei}.
\]
For $2\le i\le e$, direct binomial identities yield
\[
 \mathcal T_i(e)=e-1+i/e,\qquad
 \mathcal T_i(u+1)-\mathcal T_i(u)
       =i\frac{\binom ui}{\binom ei}+i/e.
\]
The ratio in the increment increases with $i$, since its adjacent ratio
is $(u-i)/(e-i)\ge1$ for $i<e$.  Thus
$\mathcal T_i(u)\ge\mathcal T_2(u)$.
If $P\le B$ and $\widehat R\le B-u$, \eqref{eq:width-rows} bounds
$\beta_i$ by $U_i(B)$.  It follows that
\begin{equation}\label{eq:width-propagation}
 U_2(B)\le2\binom u3
 \quad\Longrightarrow\quad
 \beta_i\le i\binom u{i+1}\quad(2\le i\le e).
\end{equation}

Now assume $e\ge6$ and $u\ge2e-2$.  Set $a_0=e-1$ and
$s=\lceil a_0/2\rceil$.  Kneser's theorem \cite{Kneser} gives
\[
 |nA|\ge\min\{u,a_0+(n-1)s\},\qquad
 W\le B(u):=u+\sum_{j\ge0}(u-a_0-sj)_+.
\]
Indeed a proper stabilizer $H$ meets $A+H$ in $c\ge2$ cosets, with
$c|H|\ge a_0$; its growth contribution $(c-1)|H|$ is at least
$\lceil a_0/2\rceil$.  This also treats composite $u$.
Because $0\in A$, the sets $nA$ increase, and
$nA=\{r:\ell_r\le n\}$.  Summing the missing residues in each layer
gives $W=u+\sum_{n\ge1}(u-|nA|)$ and the displayed bound.
Writing $\Delta=u-a_0=sk+b$, $0\le b<s$, summation gives
\[
 B(u)=u+\frac{\Delta^2}{2s}+\frac\Delta2+\frac{b(s-b)}{2s}
 \le B_*:=u+\frac{\Delta^2}{2s}+\frac\Delta2
                  +\frac{\lfloor s^2/4\rfloor}{2s},
 \qquad B(u)\le\frac{u^2}{e-1}.
\]
The last inequality follows, after clearing denominators, from
\[
 2(u^2-a_0B)=s(2s-a_0)k^2+a_0sk
                  +2(2s-a_0)kb+2b^2\ge0.
\]

The endpoint first-row bounds are
\[
 \beta_1\le\frac{e+3}{3}B-u+\frac{e(e-1)}6,
 \qquad
 \beta_1\le\frac{e+5}{4}B-u+\frac{(e-1)e(e+1)}{24}.
\]
Use the second with $B_*$ for $e=6,7,8,9$.
With $z=u-(2e-2)$, the difference from $\binom u2+1$ is, respectively,
\[
 \frac{z^2+43z+12}{24},\quad 2z+1,\quad
 \frac{3z^2+126z+211}{32},\quad
 \frac{z^2+68z+116}{16}.
\]
For $e\ge10$, the first bound with $B=u^2/(e-1)$ has difference
\[
 \frac{(e-1)(e-10)}2+1+\frac{4e-33}{6}z
                   +\frac{e-9}{6(e-1)}z^2>0.
\]
For row two the same quadratic envelope gives
\[
 2\binom u3-U_2(B(u))\ge
 \frac u6\{2u^2-(3e+6)u+6e-2\}+\binom{e-1}3.
\]
For $e=6$ the right side is $(u-10)(u-3)(u+1)/3\ge0$.
For $e=7+v$, $u=2e-2+z$, the expression in braces is
$2v^2+5vz+12v+2z^2+21z+4>0$.
Apply \eqref{eq:width-propagation} with $B(u)$ to finish every higher row.

\subsection{Five generators}

Here $q=3$.  Widths $u=4,5$ were treated above.
At $u=6,7$, \eqref{eq:width-dense-resources} gives
$(P,\widehat R)\le(8,0),(10,3)$.
The harmonic estimate $b_1\le2P/3+10/3$ gives $\beta_1\le16$ at $u=6$.
At $u=7$, $R=0$ gives $\beta_1\le20$.
If $R>0$, then $H_x\ge3$, so some surviving word has degree at least
four.  For $\tau=\sum_{n\ge3}h_n\ge1$, the degreewise harmonic bound is
\[
 b_1\le6+\frac23(P-4)-\frac16\tau<10.
\]
Integrality gives $b_1\le9$ and $\beta_1\le22$.
For rows $2$--$5$, \eqref{eq:width-rows} gives the sufficient bounds
$(44,31,8,0)$ and $(68,57,22,3)$, below the corresponding capacities
$(46,60,39,11)$ and $(77,126,119,56)$.

For $u\ge8$, Kneser's general increment two gives $4W\le u^2$.
At $u=8$ the unsimplified sum gives $W\le14$ and $U_2(14)=104<112$.
For $u=9+z$, $B=u^2/4$ gives
\[
 2\binom u3-U_2(B)=\frac{2z^3+33z^2+136z+33}{6}>0.
\]
Thus \eqref{eq:width-propagation} treats every row $i\ge2$.

For the first row, an increment smaller than three forces a stabilizer
of order two and exactly two cosets.  Consequently either all proper
increments are at least three, or the original offsets are
\[
 u=2h,\qquad \{a_1,a_2,a_3\}=\{t,h,t+h\},\qquad 0<t<h.
\]
In the first case
$6W\le u^2+u+6$.  The fourth harmonic bound $b_1\le P/2+5$ gives
$\beta_1\le B=(5u^2-7u+90)/12$.
If $A_1=\binom u2+1$, then
$A_1+1-B=((u-8)^2+17(u-8)+6)/12>0$;
integer-valuedness gives $\beta_1\le A_1$.

In the second case denote the three internal variables by $y,z,w$.
The relations $yz=xw$, $z^2=xv$, $zw=vy$ isolate $z$ in $Q$.
For $0\le j\le h$, write $t+j=2K+r$, $r\in\{0,1\}$; then
\[
 y^{h-j}w^j=x^{h-r-K}z^rv^K.
\]
The right side contains an endpoint.  Thus $Q$ is spanned by $z$ and
$y^{n-j}w^j$ with $0\le j\le n<h$, giving
$P\le\binom{h+1}2+1$.
Together with $W\le h^2$, $R\le h^2-2h$, and the third harmonic
bound, this gives
$\beta_1\le B=(11h^2-7h+30)/6$.
For $A_1=\binom{2h}2+1$,
\[
 A_1+1-B=\frac{(h-4)^2+9(h-4)+2}{6}>0.
\]
This proves the required bound for all $h\ge4$ using only a spanning
set of candidate words.

\subsection{Four generators: standard words and higher rows}

Write the two internal offsets as $0<a<b<u$, and their variables as
$y,z$.  Two internal words with the same degree and offset have the
same exponents.  Hence $Q$ is a monomial quotient of $k[y,z]$;
let $D\subset\N^2$ be its finite downset of standard exponents.
Write $\sigma=\dim_k(0:_Q(y,z))$.  Its maximal points form a socle
basis, so its Betti vector is $(1,\sigma+1,\sigma)$.
The endpoint estimates become
\begin{equation}\label{eq:width-four-rows}
 (\beta_1,\beta_2,\beta_3,\beta_4)
 \le(P+\sigma+1+R,\ 2P+\sigma+3R,\ P+3R,\ R).
\end{equation}
Kneser applied to $\{0,a,b\}$ gives $|nA|\ge\min(u,2n+1)$, whence
\[
 P\le W\le1+\lfloor u^2/4\rfloor.
\]
Ordering the $\sigma$ maximal points by increasing first coordinate
forces decreasing second coordinates.  The $j$th point dominates
$(j,\sigma-1-j)$; therefore the whole triangle
$\{(s,t):s+t\le\sigma-1\}$ lies in $D$.
In particular $\binom{\sigma+1}2\le P$ and $\sigma\le u-1$.

A surviving word gives a zero-sum-free residue sequence: a zero-sum
subword can be replaced by endpoint powers.  The long-sequence theorem
of Savchev--Chen \cite{SavchevChen} implies
\begin{equation}\label{eq:width-word-length}
 st>0\Longrightarrow s+t\le\lfloor u/2\rfloor,
 \qquad (s,t)\in D\Longrightarrow s+t\le u-2.
\end{equation}
To see the first assertion, a hypothetical longer mixed word can be
unit-normalized to labels $1,c$, $c\ge2$, with multiplicities $R_0,T_0$
and $R_0+cT_0\le u-1$.  Returning to the original labels gives
$c\alpha=\beta+Ku$, $0\le K\le c-1$, hence
$Y_\alpha^c=x^{c-1-K}Y_\beta v^K$ and $R_0<c$.
Thus $2(R_0+T_0)\le u$, a contradiction.
A pure word of length $u-1$ similarly normalizes its label to one;
the other label gives a forbidden power of degree between $2$ and
$u-1$.  This proves the second assertion.  The unit acts only on this
finite residue sequence.  It follows that $H_x\le u-3$, and hence
$R=0$ for $u\le5$.

From \eqref{eq:width-four-rows} we obtain the rational envelopes
\[
 U_2=5u^2/4-2u+4,\qquad U_3=u^2-3u+4,
 \qquad U_4=u^2/4-u+1.
\]
Put $T_i=i\binom u{i+1}$.  For $u=6+z$,
\begin{equation}\label{eq:width-four-cubic}
 12(T_2-U_2)=4z^3+45z^2+140z+36>0.
\end{equation}
Moreover $U_2-U_3=u^2/4+u$, $U_2-5U_4=3u-1$, and
\[
 8T_3=3(u-3)T_2,\qquad
 10T_4=(u-3)(u-4)T_2.
\]
For $u\ge6$ these imply $U_3\le T_3$ and $U_4\le T_4$.
At $u=5$, the bounds $(P,\sigma,R)\le(7,4,0)$ give higher rows
$(18,7,0)$, below $(25,25,9)$.
At $u=4$, $(P,\sigma,R)\le(5,2,0)$ gives $(12,5,0)$, below
$(12,7,1)$.  Width $u=3$ is a full arithmetic progression.

\subsection{Four generators: the first row}

Three relations in $C$ are
\[
 y^b=x^{b-a}z^a,\qquad
 z^{u-a}=y^{u-b}v^{b-a},\qquad
 y^{u-b}z^a=x^{u-b}v^a.
\]
Set $h=\lfloor u/2\rfloor$ and $\kappa=u-b+a$.
The axes contribute at most $2u-\kappa-1$ standard words.
The mixed words lie in the triangle of size $\binom h2$ in
\eqref{eq:width-word-length}; the third relation removes
\[
 E_h(\kappa)=
 \begin{cases}\binom{h-\kappa+2}2,&\kappa\le h,\\0,&\kappa>h.\end{cases}
\]
Consequently
\[
 P\le2u-\kappa-1+\binom h2-E_h(\kappa)
       \le2u-h-2+\binom h2,\qquad
 \sigma\le\min(h+1,\kappa).
\]
For the last bound, distinct mixed socle points have distinct first
coordinates in $1,\ldots,h-1$, while the socle triangle cannot contain
the forbidden point $(u-b,a)$.
It follows that
\[
 \beta_1\le\binom h2+u+1+\lfloor u^2/4\rfloor.
\]
Subtracting this bound from $\binom u2+1$ gives $h(h-5)/2$ for $u=2h$ and
$(h^2-3h-2)/2$ for $u=2h+1$.
This proves the comparison for $u\ge9$.
For $u=5$, $P\le7$, $\sigma\le3$, $R=0$ give $\beta_1\le11$.

At $u=6$, $H_x>2$ forces a pure word of degree four and offset at
most six.  A mixed word would violate \eqref{eq:width-word-length},
and $z^4$ has offset at least eight.  The remaining word $y^4$ forces
$a=1$; the forbidden power $y^b$ then forces $b=5$.
Thus $(a,b)=(1,5)$; there $yz=xv$ and the pure
fifth-power relations give $(P,\sigma,R)\le(9,2,4)$ and
$\beta_1\le16$.  Otherwise $R=0$ and $\beta_1\le15$.
For $u=7$, the displayed counting bounds give $\beta_1\le22$ unless
$\kappa=4$; the remaining pairs are $(1,4),(2,5),(3,6)$.
They satisfy, respectively, $z^2=yv$, $yz=xv$, $y^2=xz$, which give
$(P,\sigma)\le(7,2),(9,2),(7,2)$.
With $R\le6$ these also suffice.
For $u=8$ only $\kappa=4,5,6$ remain after counting.  The necessary
strip bounds are
\[
\begin{array}{c|c|c}
 \kappa&(a,b)&(P,\sigma)\text{ at most}\\\hline
 4&(1,5),(3,7)&(14,4)\\
 5&(1,4),(2,5),(3,6),(4,7)&(9,2)\\
 6&(1,3),(5,7)&(14,3)\\
 6&(3,5)&(9,2).
\end{array}
\]
For $\kappa=4$, the axes have at most nine points and the six-point
mixed triangle loses the forbidden point $(3,1)$ or $(1,3)$.
For $\kappa=5$, one square has an endpoint factor and the other pure
power has degree at most five; the axes and mixed words contribute at
most six and three points.  For $\kappa=6$, the outer pairs have a
forbidden cube and a forbidden seventh power, leaving at most nine
axis points and five mixed points.  In the middle pair $yz=xv$ removes
all mixed words.  The corresponding strip widths bound the number of
socle corners by two or three, as recorded in the table.
These give $\beta_1\le29$ using $R\le9$.
The lists follow by solving $b-a=u-\kappa$ and imposing
$\gcd(a,b,u)=1$.

At $u=4$, $d\ge2$ is immediate from the coarse bound $\beta_1\le8$
and $V\ge8$.  For $(a,b)=(1,2),(2,3)$ both squares vanish in $Q$.
The resulting two-by-two downset satisfies $P+\sigma\le5$, so
$\beta_1\le6$.  It remains to consider $d=1$, $(a,b)=(1,3)$.

\subsection{The remaining width-four family}

Here $m\ge6$.  Use the four homogeneous relations
\begin{equation}\label{eq:width-four-core}
 y^3-x^2z,\quad yz-xv,\quad xz^2-y^2v,\quad z^3-yv^2.
\end{equation}
Reduction decreases the internal weight $\deg_y+2\deg_z$ and leaves
the columns $x^Lv^Kb$ with $b\in\{1,y,y^2,z\}$, and $z^2v^K$.
Their degree--offset pairs are distinct, so they are independent in $C$.

Let an anchor $b$ have degree $c$ and offset $n$.
A competing free column with degree larger by $j\ge1$ has residue
$r_j=(n-jm)\bmod4$ and coordinates
\[
 K'=K-\kappa_j,\qquad L'=L-\Delta_j,\quad
 \kappa_j=\left\lceil\frac{jm-n}{4}\right\rceil,\quad
 \Delta_j=\lambda(r_j)-c-\kappa_j-j,
\]
where $(\lambda(0),\lambda(1),\lambda(2),\lambda(3))=(0,1,2,1)$.
Since $4c-n\ge0$ and $4\lambda(r)-r\le6$,
$4\Delta_j\le6-j(m+4)<0$.
Thus competition requires exactly $K\ge\kappa_j$.
A restricted column has offset equal to four times its degree minus
two, and cannot be a longer competitor when $m\ge6$.
Since $\kappa_j$ increases with $j$, the full condition reduces to
\[
 K\ge\theta_n,\qquad
 \theta_n=\max\left\{\left\lceil\frac{m-n}{4}\right\rceil,0\right\}.
\]
The defining ideal of $G$ is therefore generated by
\eqref{eq:width-four-core} and the five tails
$bv^{\theta_n}$ for $b=1,y,y^2,z,z^2$.
Indeed reduction spans the quotient by the surviving columns, whose
images are independent semigroup monomials in their exact order.

For $m=6$, two core relations and four tails suffice, giving
$\beta_1\le6<9$.
For $m>6$, write $m=4Q+\rho$.  In anchor order $1,y,y^2,z,z^2$,
the five thresholds are
\[
\begin{array}{c|ccccc}
 \rho&\theta_0&\theta_1&\theta_2&\theta_3&\theta_6\\\hline
 0&Q&Q&Q&Q&Q-1\\
 1&Q+1&Q&Q&Q&Q-1\\
 2&Q+1&Q+1&Q&Q&Q-1\\
 3&Q+1&Q+1&Q+1&Q&Q.
\end{array}
\]
A tail is redundant whenever a divisor anchor has a threshold no
larger than its own.  Retain all four core relations and, according to
$\rho$, respectively two, four, four, or two tails.  Hence
\[
\begin{array}{c|rrrr}
 \rho&0&1&2&3\\\hline
 \text{number of generators at most}&6&8&8&6\\
 F_1(m,4)&7&10&9&8.
\end{array}
\]
This completes the first row for every $m$ in the remaining family.

The preceding cases cover $e=4$ and all $u\ge3$.
For $e\ge6$, the ranges $u=e-1,e$, $e+1\le u\le2e-3$, and
$u\ge2e-2$ are covered by the smallest, dense, and sparse width
arguments; the five-generator discussion covers $e=5$.
Finally $\beta_0=1$ on both sides and $\beta_i(G)=0$ for $i>e$.
Thus every narrow source satisfies the Betti comparison with its
literal interval completion, over the original field.

\section{Small embedding dimension and conclusion}
\label{sec:synthesis}
We complete the argument at the small-generator boundary, then collect
the cases. All Betti numbers below use the original coefficient field.

\subsection{One, two, and three generators}
If $e=1$, the gcd condition forces $S=\N$. Its interval completion
is the literal semigroup $\langle1\rangle=\N$, and both Betti sequences
are $(1,0,0,\ldots)$.

If $e=2$, write $S=\langle m,n\rangle$, with $m<n$ and
$\gcd(m,n)=1$. The completed semigroup ring is
$k[[X,Y]]/(Y^m-X^n)$, and its tangent cone is
$k[X,Y]/(Y^m)$. Indeed, the least-degree form of any nonzero multiple
of $Y^m-X^n$ is the corresponding least-degree form times $Y^m$.
The source Betti sequence is therefore $(1,1,0,\ldots)$.
The target has at least two minimal generators and
$F_1(m,v)\ge1$, proving the comparison.

For $e=3$, the comparison also follows from
\cite[Theorem~1.1]{LanTuVu}. We give the one-letter version of the
endpoint argument to record the proof used in the formal development.
In the ordinary-width case Section~\ref{sec:ordinary} applies.
Otherwise normalize as in \eqref{eq:normalization}, now with $q=1$.
If $u=2$ the source is a full arithmetic sequence, covered by the
full-progression comparison in Section~\ref{sec:width}.
Suppose $u\ge3$.

Use the endpoint core and its invariants $P_C,R$ from
Section~\ref{sec:endpoint}. Every minimal internal word is a pure
power $Y^n$ with $n<u$, and its degree determines it. Hence $P_C\le u$.
For each minimal competitor of positive first coordinate, choose a
minimal core point below its predecessor and record the positive jump
$j$. The competitor is recovered from that point and $j$, so the
assignment is injective. If the core word has degree $n$, the number of
permitted jumps is at most $(n-3)_+$, where $(x)_+=\max\{x,0\}$.
Competitors of first coordinate zero contribute at most one per
residue and cancel the one-per-residue term in the socle count.
Consequently,
\[
 R\le\sum_{n=0}^{u-1}(n-3)_+=\binom{u-3}{2}.
\]
The one-variable core has Betti sequence $(1,1)$; its Euler
characteristic, or its presentation as a nonzero finite-length cyclic
$k[Y]$-module, gives the positive row. The compensated endpoint bound
therefore yields
\[
 (\beta_1,\beta_2,\beta_3)(G_S;k)
 \le (P_C+R+1,\;P_C+2R,\;R)
\]
componentwise. The first two rows are paid for by
\[
 u+\binom{u-3}{2}\le\binom u2,\qquad
 \binom{u-3}{2}\le\binom u3
\]
and \eqref{eq:low-capacities}. For the third row, $R=0$ if $u=3,4$;
if $u\ge5$, use $\binom{u-3}{2}\le\binom u2$ and the third-row
capacity. Monotonicity transfers these normalized bounds to the
physical width $du$.

\subsection{Exhaustion of the parameter space}
For $e\ge4$, the integer width satisfies exactly one of
$w\ge m-1$ and $w\le m-2$. The first case is proved in
Section~\ref{sec:ordinary}. In the second, the normalization retains
the same semigroup and the same multiplicity. The narrow-width proof
in Section~\ref{sec:width} applies for $e\ge5$ and for $e=4$ by
its two-internal-variable analysis. Thus the cases are disjoint and
exhaustive:
\[
\begin{array}{c|l}
\text{source parameters}&\text{argument}\\\hline
e=1,2,3&\text{small-generator comparison above}\\
e\ge4,\;w\ge m-1&\text{ordinary completion}\\
e=4,\;w\le m-2&\text{two internal variables}\\
e\ge5,\;w\le m-2&\text{general narrow-width comparison}.
\end{array}
\]
Each case proves all potentially nonzero positive source rows, including
the row $i=e$. The zeroth Betti number is one on both sides. For $i>e$
the source Koszul complex is zero, so the comparison is automatic,
regardless of the target's embedding dimension. Completion and any
temporary scalar extension are removed by Proposition~\ref{prop:bridge}.
This proves Theorem~\ref{thm:main}.

\begin{remark}
The two resolutions can have different lengths. For example,
$S=\langle6,7,9,11\rangle$ has four minimal generators, so
$\beta_5(G_S;k)=0$. Its completion is $\langle6,7,8,9,10,11\rangle$,
whose Betti sequence is
$(1,15,40,45,24,5)$ by \eqref{eq:ordinary-capacity}.
Thus the comparison continues beyond the source's last possible row;
one does not identify the source and target polynomial rings.
\end{remark}

\appendix
\section{Formal Verification and Reproducibility}
\label{sec:formalization}
The formal theorem is
\begin{center}
\texttt{NSP14.NumericalSemigroup.}\\
\texttt{p14\_totalBetti\_le\_intervalCompletion}.
\end{center}
Its parameters are a numerical semigroup $S$, an arbitrary field $k$,
and a natural number $i$. The source and target are the associated
graded algebras of the actual supported formal-power-series semigroup
rings. The presentation variables are their respective minimal atoms.
The invariant \texttt{totalBetti} is defined as the dimension of the
full-variable Koszul homology, and a proved isomorphism identifies this
homology with ordinary $\Tor$ over that polynomial ring.

The source used for this manuscript is the repository
\href{https://github.com/stksk310/NS-P14-Lean}{\texttt{stksk310/NS-P14-Lean}},
at commit
\begin{center}\small
\texttt{a5e122a4d09536b1609c5faa8d419f8faf81671f}.
\end{center}
The compiler is Lean \texttt{v4.36.0-rc2}. The mathlib revision is
\begin{center}\small
\texttt{139105d28c7685a154cb3093f79edb2e16708099}.
\end{center}
The principal interfaces are organized as follows.
\begin{center}\small
\begin{tabular}{p{.28\linewidth}p{.62\linewidth}}
\toprule
Mathematical role&Lean module\\\midrule
Completed-ring bridge&\raggedright\arraybackslash
\texttt{Semigroup.FormalTangentAlgebra},
\texttt{Homological.ActualCoefficientBaseChange}\\
Exact target&\texttt{ArithmeticSequence.BettiFormula},
\texttt{Main.IntervalTarget}\\
Ordinary comparison&\texttt{Ordinary.Closure}\\
Endpoint mechanism&\texttt{Endpoint} and its homological consumers\\
Narrow-width comparison&\texttt{Width.Closure}, \texttt{NarrowFour.Comparison}\\
Small generators&\texttt{SmallEmbedding.Base},
\texttt{SmallEmbedding.ThreeComparison}\\
Exhaustion and assembly&\texttt{Main.Partition}, \texttt{Main.Branches},
\texttt{Main}\\\bottomrule
\end{tabular}
\end{center}
All module names have prefix \texttt{NSP14}.
The released completion record is
\path{verification/run20/RUN20_COMPLETE_20261116_R03.json}; it binds
the proof snapshot, validation evidence, and a separate AI-assisted audit.
The recorded transitive axioms are Lean's standard
\texttt{propext}, \texttt{Classical.choice}, and \texttt{Quot.sound}.
Published mathematical inputs used in the proof of
Theorem~\ref{thm:main} are proved, in the forms required there, in the
formal development. Corollary~\ref{cor:width-bound} additionally uses
the standard filtered-Koszul comparison with the complete local ring.
The repository's original fifty-entry dependency index and its final
paper-correspondence ledger give finer declaration-level references.

Lean kernel checking is distinct from the generative-AI audit workflow:
the formal theorem is checked against its definitions and proof terms.
The accompanying Python scripts orchestrate compilation, declaration
and axiom checks, and the preservation of verification evidence.
Finite arithmetic checks supplement the general proofs and are not
used to extrapolate an unbounded assertion from sampled cases.

To reproduce the released verification, retain the repository's Git
history and check out the proof snapshot identified by the completion
record, rather than the later commit that packages the release:
\begin{center}\small
\texttt{1a5c6bac9a4fd38a8bc5d13b4c7372e654466e6e}.
\end{center}
Install the pinned compiler and dependencies, including their build
caches. With the pinned Lean executables on \texttt{PATH}, run
\begin{small}
\begin{verbatim}
python3 scripts/verify_run20.py \
  --fresh-project-cache \
  --output-dir verification/run20/runs/independent_new_id
\end{verbatim}
\end{small}
Choose a new output directory for each run. The verifier rebuilds the
project library, intermediate code, tests, and audits while retaining
the fixed dependency cache. The audit archive does not contain an
offline compiler and dependency installation. The archived release
evidence and a new execution of these instructions are distinct records.

\section*{Data Access Statement}
The Lean sources, pinned dependency manifest, verification scripts, and
archived verification evidence supporting the tangent-cone comparison
are publicly available in the versioned repository
\begin{center}
\href{https://github.com/stksk310/NS-P14-Lean/tree/a5e122a4d09536b1609c5faa8d419f8faf81671f}{\texttt{stksk310/NS-P14-Lean}}.
\end{center}
The release commit is
\begin{center}\small
\texttt{a5e122a4d09536b1609c5faa8d419f8faf81671f}.
\end{center}
Appendix~\ref{sec:formalization} identifies the proof snapshot and
reproduction procedure. The release includes the completion record and
the archive \path{artifacts/NS_P14_LEAN_RUN20_COMPLETE_20261116_R03.zip}.

\section*{Acknowledgements}
\textit{Funding.} This research received no external funding.

\textit{Conflicts of interest.} The author declares no conflicts of interest.

\bibliographystyle{amsplain}
\bibliography{references}
\end{document}
